\documentclass[11pt,a4paper]{article}
\input{../../sablony/pripravy/priprava.sty}
\usepackage{tikz,pgfplots}
\pgfplotsset{compat=1.18}
\nastavPripravu{1. ročník}{Kvadratické rovnice a nerovnice}{126}{Soustavy kvadratických nerovnic}
\begin{document}
\blok{Soustavy kvadratických nerovnic}
Průnik dvou nebo více intervalových řešení; opakování množin.
\blok{Metoda a návaznost na graf}
\begin{itemize}
\item Nejprve vyřeš nerovnice zvlášť a potom najdi průnik jejich řešení.
\item Nakresli každou množinu na číselnou osu; prázdný průnik je legitimní výsledek.
\end{itemize}
\blok{Proč průnik?}
Nerovnice popisují dvě podmínky, které mají platit současně. Nejprve označíme například
$K_1=\langle-3;3\rangle$ a $K_2=(1;5)$, pak na stejné číselné ose vyznačíme jejich průnik
$K_1\cap K_2=(1;3\rangle$. Nulové body se započítají podle ostrého či neostrého znaménka.
\blok{Společně řešené příklady}
\begin{prikladReseniBox}[title={Příklad 1}]
$\begin{cases}x^{2} - 9<0\\x^{2} - 4x + 3>0\end{cases}$
\tcblower
\[
\begin{aligned}
K_1&=(-3;3)\\[2pt]
K_2&=(-\infty;1)\cup(3;+\infty)\\[2pt]
K&=K_1\cap K_2=(-3;1)
\end{aligned}
\]
\end{prikladReseniBox}
\begin{prikladReseniBox}[title={Příklad 2}]
$\begin{cases}x^{2} - 5x + 6\le0\\x^{2} - 16\ge0\end{cases}$
\tcblower
\[
\begin{aligned}
K_1&=\langle2;3\rangle\\[2pt]
K_2&=(-\infty;-4\rangle\cup\langle4;+\infty)\\[2pt]
K&=K_1\cap K_2=\varnothing
\end{aligned}
\]
\end{prikladReseniBox}
\blok{Samostatné procvičování}
\begin{zadaniReseniBox}[title={Samostatná práce 1}]
$\begin{cases}x^{2} - 4\ge0\\x^{2} - 6x + 5<0\end{cases}$
\tcblower
\[
\begin{aligned}
K_1&=(-\infty;-2\rangle\cup\langle2;+\infty)\\[2pt]
K_2&=(1;5)\\[2pt]
K&=K_1\cap K_2=\langle2;5)
\end{aligned}
\]
\end{zadaniReseniBox}
\blok{Volitelná rezerva}
\poznamka{Studenti si vymyslí obdobný příklad se změnou jednoho koeficientu a vysvětlí, jak se změní počet kořenů či řešení.}
\blok{Domácí úkol}
\begin{zadaniReseniBox}[title={Domácí úkol}]
\begin{enumerate}[label=\alph*),itemsep=2pt]
\item $\begin{cases}x^{2} - 25<0\\x^{2} - 4x + 3>0\end{cases}$
\item $\begin{cases}x^{2} - 7x + 10\le0\\x^{2} - 9\ge0\end{cases}$
\item $\begin{cases}x^{2} - 1>0\\x^{2} - 2x - 3<0\end{cases}$
\item $\begin{cases}x^{2} - 5x + 6>0\\x^{2} - 16<0\end{cases}$
\item $\begin{cases}x^{2} - 4x + 4\le0\\x^{2} - 16\ge0\end{cases}$
\item $\begin{cases}x^{2} - 9>0\\x^{2} - 8x + 12<0\end{cases}$
\end{enumerate}
\tcblower
\begin{enumerate}[label=\alph*),itemsep=2pt]
\item
\[\begin{aligned}
K_1&=(-5;5)\\[2pt]
K_2&=(-\infty;1)\cup(3;+\infty)\\[2pt]
K&=K_1\cap K_2=(-5;1)\cup(3;5)
\end{aligned}\]
\item
\[\begin{aligned}
K_1&=\langle2;5\rangle\\[2pt]
K_2&=(-\infty;-3\rangle\cup\langle3;+\infty)\\[2pt]
K&=K_1\cap K_2=\langle3;5\rangle
\end{aligned}\]
\item
\[\begin{aligned}
K_1&=(-\infty;-1)\cup(1;+\infty)\\[2pt]
K_2&=(-1;3)\\[2pt]
K&=K_1\cap K_2=(1;3)
\end{aligned}\]
\item
\[\begin{aligned}
K_1&=(-\infty;2)\cup(3;+\infty)\\[2pt]
K_2&=(-4;4)\\[2pt]
K&=K_1\cap K_2=(-4;2)\cup(3;4)
\end{aligned}\]
\item
\[\begin{aligned}
K_1&=\{2\}\\[2pt]
K_2&=(-\infty;-4\rangle\cup\langle4;+\infty)\\[2pt]
K&=K_1\cap K_2=\varnothing
\end{aligned}\]
\item
\[\begin{aligned}
K_1&=(-\infty;-3)\cup(3;+\infty)\\[2pt]
K_2&=(2;6)\\[2pt]
K&=K_1\cap K_2=(3;6)
\end{aligned}\]
\end{enumerate}
\end{zadaniReseniBox}
\poznamka{Orientační čas: úvod 5 min, společné příklady 18 min, samostatně 15 min, kontrola a shrnutí 7 min. Rezerva je dobrovolná. Číselné parametry modelových aplikací nejsou skutečně naměřená data.}
\end{document}
